% ===========================================================================
\section{Related Work}
\label{sec:related}
% ===========================================================================

\subsection{Slow Fault Testing}

\para{Xinda} Ganesan~\etal{}~\cite{xinda} (NSDI~'25) introduced systematic
slow-fault testing through exhaustive grid exploration. \xinda{} provides a
comprehensive pipeline for deploying systems, injecting faults at configured
severities, and reporting danger zones. \system{} builds on \xinda{}'s insight
that slow faults are a critical yet under-tested failure class, but replaces
exhaustive enumeration with adaptive binary search to achieve minimization with
fewer trials. \system{} is complementary to \xinda{}: operators can use
\xinda{}'s grid for initial exploration and \system{}'s minimizer for precise
boundary characterization.

\para{Anduril} Yao~\etal{}~\cite{anduril} (ATC~'25) targets \emph{crash
and exception} bugs using feedback-directed fault injection. \anduril{} builds
a causal DAG from runtime traces and preferentially injects faults at high-impact
points. Unlike \system{}, \anduril{} requires source-level instrumentation,
targets discrete event triggers rather than sustained degradation, and does not
produce minimal severity characterizations.

\para{PRE-FAIL} Joshi~\etal{}~\cite{prefail} proposed programmable failure
injection at the I/O boundary. PRE-FAIL intercepts filesystem and network calls
to inject configured failures, enabling precise control over failure timing.
However, PRE-FAIL requires application modification (linking against its library)
and does not include minimization or systematic exploration.

\subsection{Fault Injection Frameworks}

\para{Jepsen} Kingsbury~\cite{jepsen} provides a comprehensive framework
for testing distributed systems' consistency guarantees under network
partitions and node failures. Jepsen focuses on \emph{correctness} (safety
violations) rather than \emph{performance} (liveness violations under
degradation). \system{} targets the complementary problem of characterizing
\emph{how much} degradation triggers visible failure, not whether the system
maintains consistency.

\para{Chaos Monkey and Chaos Engineering} Netflix's Chaos Monkey~\cite{chaos-monkey}
and the broader chaos engineering movement~\cite{chaos-engineering-book}
inject failures in production to validate resilience. These tools verify
that systems \emph{survive} random failures but do not characterize minimal
triggers or produce boundary specifications. \system{}'s minimal recipes
could inform chaos engineering experiments by targeting the most critical
severity range.

\para{Gremlin} Commercial fault injection platform providing network, resource,
and state faults~\cite{gremlin}. Gremlin enables scenario-based testing but
relies on manual scenario design. \system{}'s automated minimization
complements manual scenario tools by discovering the minimal configurations
that actually trigger failures.

\para{LitmusChaos} Cloud-native chaos engineering framework for Kubernetes
environments~\cite{litmuschaos}. Provides a library of pre-built chaos
experiments (pod kill, network loss, disk fill) with GitOps-based workflow.
Like other chaos tools, focuses on binary pass/fail testing rather than
boundary characterization.

\subsection{Lineage-Driven Fault Injection}

\para{Molly/LDFI} Alvaro~\etal{}~\cite{alvaro-molly,alvaro-ldfi} introduced
lineage-driven fault injection, which uses data provenance to identify faults
that may violate program outcomes. LDFI works backward from a successful
outcome to find minimal fault combinations that break it. While conceptually
similar to minimization, LDFI operates on discrete message-level faults
(dropped/delayed messages) rather than continuous severity parameters, and
requires a formal specification of program correctness. \system{}'s
symptom-guided oracle is a lightweight alternative to formal specifications.

\para{Filibuster} Meiklejohn~\etal{}~\cite{filibuster} adapt LDFI to
microservice architectures, injecting RPC-level failures guided by service
dependency analysis. Filibuster targets inter-service communication failures
(timeouts, HTTP errors) rather than intra-cluster consensus degradation.
The two approaches are complementary: Filibuster for service mesh resilience,
\system{} for individual system fault boundaries.

\subsection{Delta Debugging and Minimization}

\para{Delta debugging} Zeller and Hildebrandt~\cite{zeller-delta-debugging}
introduced automated test case minimization for crash-reproducing inputs.
Given a failing test input, delta debugging systematically removes portions
until a minimal reproducing input remains. \system{}'s greedy dimensional
minimizer is conceptually related but operates over continuous fault parameters
rather than discrete input elements, and uses binary search rather than
ddmin's divide-and-partition strategy.

\para{C-Reduce} Regehr~\etal{}~\cite{creduce} extended delta debugging with
transformation-based reduction for compiler bug test cases. C-Reduce's
multi-pass strategy (each pass applies a different reduction transformation)
inspired \system{}'s phase-based approach (each phase reduces a different fault
dimension).

\para{Test case prioritization} Rothermel~\etal{}~\cite{rothermel-tcp}
address the problem of ordering test cases for early fault detection. While
\system{} does not prioritize existing tests, its binary search strategy
implicitly front-loads the most informative trials (mid-range values that
maximally divide the search space).

\subsection{Gray Failure and Partial Failures}

\para{Gray failure} Huang~\etal{}~\cite{huang-gray-failure} characterized
``gray failures'' where system components enter degraded states undetected
by health checks. Their taxonomy (observation gap, reaction gap) directly
motivates \system{}: our minimal recipes quantify exactly where the
observation gap begins (the severity below which degradation is undetectable
by current monitoring).

\para{Fail-slow hardware} Gunawi~\etal{}~\cite{gunawi-cbs} cataloged
``fail-slow'' hardware incidents in production cloud systems, finding that
storage and network devices frequently enter degraded-performance states.
Their taxonomy informs our fault model: network, filesystem, CPU, and
memory faults correspond to the most common fail-slow categories.

\para{Cloud reliability studies} Lu~\etal{}~\cite{lu-cloud-reliability}
analyzed 597 cloud system failures and found that 33\% involved
performance-related issues rather than hard crashes. Their finding that
``performance bugs are harder to detect, diagnose, and fix'' motivates
\system{}'s focus on making slow-fault boundaries explicit and testable.

\subsection{Distributed System Testing}

Model checking approaches (SAMC/FlyMC~\cite{leesatapornwongsa-samc}) and
static analysis (CrashTuner~\cite{crashtuner}) target correctness bugs with
discrete triggers rather than continuous performance degradation boundaries.
